% Introduction

%Context
Free/Libre/Open Source Software (FLOSS) projects
have been used in Software Engineering Education (SEE)
to address the need for more realistic settings
that could reduce the gap between
software engineering courses and industry needs.
%Several aspects of open source development have pedagogical benefit in a computer science or software engineering program. 
The study of FLOSS projects exposes students to 
large projects, best practices and professional tools 
for version control, build automation, and project management, 
while allowing them to experience the collaborative, social dynamics of team working on these large projects~\cite{Horstmann:2009}.

The most recent systematic mapping study (SMS)  \cite{Petersen,Budgen}
that investigated the use of FLOSS projects in SEE,
structured the research area according to nine facets and
provided an overview of initiatives and research gaps
based on work published prior to 2013 \cite{2015:CSE:nascimento}.
% Main Results 
This SMS identified active research communities, long-term projects and venues for publication, and provided an overview of objectives, contributions and methods applied in 72 studies over 14 years (1998-2012). 
This work provided a baseline for the scholarship on SEE with FLOSS, describing previous work in terms of research/scholarship type, software engineering area, curriculum choice, FLOSS choice, control level, goals and assessment. Results from selected papers showed a gap in the use of sound research methods, a lack of more active learning approaches as well a lack of educational theory foundations. Most selected papers were published in conferences, describing either solution proposals or experience reports on comprehensive software engineering courses based on the traditional project method approach. Interestingly, this paper showed both the community interest of using FLOSS in SEE and the need to improve scholarship to turn this area into a mature field of research.

\begin{figure*}[tb]
\centering
\includegraphics[width=1\linewidth]{fig/update-process.jpg}
\caption{The update process, adapted from \cite{Petersen,2015:CSE:nascimento}.}
\label{fig_smsProcess3}
\end{figure*}
 
The SIGCSE 2016 Panel on ``How to Use Open Source Software in Education'' capitalized on the uses of FLOSS in education, 
pointing out to barriers to its use in educational settings and 
the need for a debate to better take advantage of the collaborative ways to solve hard problems that open source communities and their ecosystems bring \cite{Bishop:2016}.

% The need for an update. WHY? Justify.
Therefore, an updated SMS represents an important 
asset in a new research field that brings together 
FLOSS and SEE, where the number of studies and the maturity of empirical methods used are increasing.

% Talk a little bit about the need for an update.
% This work:
In this work, we present an update of the SMS on 
the use of FLOSS projects in 
SEE~\cite{2015:CSE:nascimento}.
% Talk a little bit more about the update.
We retrieved and analyzed papers published from 2013 to 2017,
and classified them according to the facets used in the original SMS.
% How many facets are presented here and WHY?
We analyzed the new results and compared them
with those from the previous SMS 
to confirm trends and discover new ones.

Our main findings point out both a confirmation of previous research trends as well as some changes that slowly move the field into a more mature research topic. Confirmed trends are: 
(1)~the use of FLOSS to learn comprehensively on software engineering; (2)~curriculum choice particularly focused on regular courses; and 
(3)~assessment mainly based on reports and software artifacts. 
On the other hand, new trends point to: 
(4)~increased use of experience reports and a reduction in solution proposals; 
(5)~growing interest of using FLOSS in software maintenance and evolution; 
(6)~increased use of approaches where instructors have no control, and students interact with the FLOSS community in real tasks; 
(7)~less use of predefined FLOSS projects and more free project choice by students; 
(8)~balance between goals of learning FLOSS and learning SE concepts; and (9)~the use of automated test suites as a relevant assessment type.    

%\textit{Paper Structure.}
The rest of the paper is organized as follows.
Section~\ref{sec:sms:process} presents information about process steps 
performed and outcomes generated in the context of the SMS update.
Section~\ref{sec:updated:sms} presents the results of the updated mapping and compares them with results from the 2013 SMS.
%Section~\ref{sec:trends} discusses ... and trends.
Section~\ref{sec:conclusion} presents contributions and discloses opportunities for future research.


